%% ============================================================
%% Template Elsevier — elsarticle
%% Instrucciones:
%%   - Compila con: pdflatex main.tex → bibtex main → pdflatex x2
%%   - O usa Overleaf: sube la carpeta completa como proyecto
%%   - Sustituye los bloques marcados con TODO
%% ============================================================

\documentclass[preprint,12pt,authoryear]{elsarticle}

%% --- Paquetes básicos ----------------------------------------
\usepackage[utf8]{inputenc}
\usepackage[T1]{fontenc}
\usepackage[spanish,english]{babel}   % cambia a english si el paper es en inglés
\usepackage{amsmath, amssymb}
\usepackage{graphicx}
\usepackage{booktabs}                 % tablas de calidad
\usepackage{hyperref}
\usepackage{cleveref}
\usepackage{siunitx}                  % unidades (°C, m/s, ppm...)
\usepackage{listings}                 % bloques de código
\usepackage{xcolor}
\usepackage{microtype}

%% --- Configuración de listings (código) ---------------------
\lstset{
  basicstyle=\ttfamily\small,
  breaklines=true,
  frame=single,
  numbers=left,
  numberstyle=\tiny,
  keywordstyle=\color{blue},
  commentstyle=\color{gray},
  stringstyle=\color{red!70!black}
}

%% --- Información del journal (cambia según la revista) ------
\journal{Journal Name}   % TODO: nombre de la revista

%% ============================================================
\begin{document}

\begin{frontmatter}

%% --- Título -------------------------------------------------
\title{%
  % TODO: título del paper (claro, específico, con palabras clave)
  A Data-Driven Approach for Monitoring and Anomaly Detection
  in Greenhouse Sensor Networks
}

%% --- Autores ------------------------------------------------
% TODO: añade o quita autores según corresponda
\author[inst1]{Nombre Apellido\corref{cor1}}
\ead{autor@universidad.es}

\author[inst1,inst2]{Segundo Autor}
\ead{segundo@universidad.es}

\cortext[cor1]{Corresponding author}

\affiliation[inst1]{
  organization={Nombre del Grupo / Departamento},
  addressline={Calle o Campus},
  city={Ciudad},
  postcode={Código postal},
  country={Spain}
}

\affiliation[inst2]{
  organization={Segunda institución (si aplica)},
  city={Ciudad},
  country={Country}
}

%% --- Abstract -----------------------------------------------
\begin{abstract}
% TODO: 150–250 palabras. Estructura recomendada:
%   (1) Contexto y problema  (2) Objetivo  (3) Método
%   (4) Resultados principales  (5) Conclusión/impacto

This paper presents a framework for real-time monitoring and
anomaly detection in greenhouse environments using low-cost
sensor networks. We address the problem of \ldots
The proposed approach combines \ldots
Experimental results on a seven-day synthetic dataset show
that the method achieves \ldots
The results suggest that \ldots
\end{abstract}

%% --- Palabras clave -----------------------------------------
\begin{keyword}
% TODO: 4–6 palabras clave separadas por \sep
greenhouse monitoring \sep
anomaly detection \sep
sensor data \sep
data-driven methods \sep
Python \sep
reproducible research
\end{keyword}

\end{frontmatter}

%% ============================================================
\section{Introduction}
\label{sec:introduction}

% TODO: ~0.5–1 página
% Estructura: contexto general → problema concreto →
%             limitaciones del estado del arte →
%             contribución de este trabajo → estructura del paper

Greenhouse automation and precision agriculture have attracted
increasing attention in recent years due to \ldots
~\citep{Ref1}.

Despite advances in sensor technology, several challenges remain:
(i)~\ldots; (ii)~\ldots; (iii)~\ldots.

The main contributions of this work are:
\begin{itemize}
  \item A synthetic benchmark dataset representing seven days of
        multivariate sensor readings, including two simulated fault
        events (\cref{sec:data}).
  \item A reproducible analysis pipeline implemented in Python and
        distributed as a Docker image (\cref{sec:methods}).
  \item Quantitative evaluation of \ldots (\cref{sec:results}).
\end{itemize}

The remainder of the paper is organized as follows.
\cref{sec:related} reviews related work.
\cref{sec:data} describes the dataset.
\cref{sec:methods} details the proposed methodology.
\cref{sec:results} presents the experimental evaluation.
\cref{sec:conclusion} concludes the paper.

%% ============================================================
\section{Related Work}
\label{sec:related}

% TODO: ~0.5–1 página
% Agrupa trabajos por tema/enfoque, no hagas una lista cronológica.
% Acaba con un párrafo que identifique el hueco que llena este paper.

Several approaches have been proposed for greenhouse monitoring.
\citet{Ref2} proposed \ldots, while \citet{Ref3} focused on \ldots
In contrast to these works, our approach \ldots

%% ============================================================
\section{Dataset and Experimental Setup}
\label{sec:data}

% TODO: describe los datos, el entorno y las condiciones del experimento

\subsection{Sensor Variables}

The dataset covers seven days of measurements sampled every
\SI{15}{\minute} (\num{672} samples per variable).
\cref{tab:variables} summarizes the monitored variables.

\begin{table}[ht]
\centering
\caption{Sensor variables recorded in the dataset.}
\label{tab:variables}
\begin{tabular}{llll}
\toprule
Variable & Symbol & Unit & Range \\
\midrule
Indoor temperature  & $T_\text{int}$ & \si{\celsius}   & 10–35 \\
Outdoor temperature & $T_\text{ext}$ & \si{\celsius}   & 5–28  \\
Indoor humidity     & $H_\text{int}$ & \si{\percent}   & 45–95 \\
Outdoor humidity    & $H_\text{ext}$ & \si{\percent}   & 30–85 \\
CO\textsubscript{2} concentration & $c_\text{CO2}$ & ppm & 400–1500 \\
Wind speed          & $v_\text{wind}$& \si{\metre\per\second} & 0–12 \\
\bottomrule
\end{tabular}
\end{table}

\subsection{Simulated Fault Events}

Two fault events were injected into the dataset:
\begin{enumerate}
  \item \textbf{Heating failure} (Day~3): indoor temperature drops
        \SI{8}{\celsius} below the expected profile over \SI{6}{\hour}.
  \item \textbf{Ventilation failure} (Day~5): CO\textsubscript{2}
        concentration rises by up to \SI{700}{ppm} over \SI{4}{\hour}.
\end{enumerate}

%% ============================================================
\section{Methodology}
\label{sec:methods}

% TODO: describe el método propuesto con suficiente detalle para
%       que sea reproducible. Usa subfiguras si es necesario.

\subsection{Overview}

\cref{fig:pipeline} illustrates the proposed processing pipeline.
% TODO: añade la figura en paper/figuras/pipeline.pdf o .png
% \begin{figure}[ht]
%   \centering
%   \includegraphics[width=0.85\linewidth]{figuras/pipeline}
%   \caption{Overview of the proposed pipeline.}
%   \label{fig:pipeline}
% \end{figure}

\subsection{Preprocessing}

Raw sensor readings are normalized using min-max scaling:
\begin{equation}
  \hat{x} = \frac{x - x_{\min}}{x_{\max} - x_{\min}}
  \label{eq:normalization}
\end{equation}

\subsection{Anomaly Detection}

% TODO: describe el detector, umbral, ventana temporal, etc.

%% ============================================================
\section{Results and Discussion}
\label{sec:results}

% TODO: presenta resultados cuantitativos con tablas y figuras.
%       Interpreta los números; no los enumeres sin explicación.

\cref{tab:results} summarizes the detection performance.

\begin{table}[ht]
\centering
\caption{Anomaly detection performance.}
\label{tab:results}
\begin{tabular}{lccc}
\toprule
Method & Precision & Recall & F1-score \\
\midrule
Baseline    & 0.XX & 0.XX & 0.XX \\
Proposed    & \textbf{0.XX} & \textbf{0.XX} & \textbf{0.XX} \\
\bottomrule
\end{tabular}
\end{table}

% TODO: discute por qué funciona, limitaciones y trabajo futuro

%% ============================================================
\section{Conclusion}
\label{sec:conclusion}

% TODO: 1 párrafo de resumen, 1 de hallazgos principales,
%       1 de trabajo futuro. Sin introducir ideas nuevas.

This paper presented \ldots
The main findings are: (i)~\ldots; (ii)~\ldots; (iii)~\ldots
Future work will focus on \ldots

%% ============================================================
%% Referencias
\bibliographystyle{elsarticle-harv}   % autor-año (Harvard)
% Alternativa numérica: elsarticle-num
\bibliography{referencias}

\end{document}
